\documentclass[11pt,oneside]{report}

\usepackage[margin=1.1in]{geometry}
\usepackage[T1]{fontenc}
\usepackage[utf8]{inputenc}
\usepackage{lmodern}
\usepackage{booktabs}
\usepackage{longtable}
\usepackage{array}
\usepackage{listings}
\usepackage{xcolor}
\usepackage{microtype}
\usepackage{setspace}
\usepackage[hidelinks]{hyperref}

\setstretch{1.12}

\definecolor{codebg}{rgb}{0.96,0.96,0.96}
\definecolor{keyword}{rgb}{0.13,0.29,0.53}
\definecolor{comment}{rgb}{0.35,0.45,0.35}
\definecolor{strcol}{rgb}{0.55,0.25,0.25}

\lstdefinelanguage{Rust}{
  keywords={fn,let,mut,pub,impl,enum,struct,trait,match,if,else,return,use,const,for,in,as,dyn,Self,self,crate,mod,where,move,ref,true,false,Some,None,Ok,Err},
  sensitive=true,
  comment=[l]{//},
  morecomment=[s]{/*}{*/},
  string=[b]",
}

\lstset{
  basicstyle=\ttfamily\footnotesize,
  keywordstyle=\color{keyword}\bfseries,
  commentstyle=\color{comment}\itshape,
  stringstyle=\color{strcol},
  backgroundcolor=\color{codebg},
  frame=single,
  framerule=0pt,
  breaklines=true,
  columns=fullflexible,
  showstringspaces=false,
  upquote=true,
  xleftmargin=1em,
  xrightmargin=1em
}

\newcommand{\wpm}{\texttt{wpm}}
\newcommand{\projver}{v26.6.10}

\title{
  {\Huge \textbf{wasm4pm}}\\[0.6em]
  {\LARGE A Process-Mining Platform with a Lawful Cognition Layer}\\[0.8em]
  {\large Technical Monograph, Release \projver}
}
\author{Sean Chatman}
\date{June 10, 2026}

\begin{document}

\maketitle

\begin{abstract}
\noindent
\textbf{wasm4pm} is a process-mining platform built as two cooperating layers:
a Rust core compiled to WebAssembly that hosts 60 process-mining algorithms,
and a TypeScript monorepo of twelve packages whose published command-line
interface, \wpm{}, drives discovery, conformance checking, and predictive
analysis from event logs. On top of the mining core sits a \emph{cognition
layer} --- a ``periodic table'' of 52 classical symbolic-AI reasoning
\emph{breeds}, each reconstructed from its source paper, each gated by an
anti-cheat oracle calculus, and each forced through a lawful runtime
lifecycle in which every successful run must yield an inference trace, pass
postconditions, and replay against an object-centric event-log (OCEL)
lifecycle model with conformance fitness exactly $1.0$. The platform's
operating doctrine, after van der Aalst, is \emph{process truth}: if the code
says it worked but the event log cannot prove a lawful process happened, then
it did not work. This monograph documents the entire system at release
\projver{}: architecture, the cognition periodic table and its ten-rung
standing ladder, the universal anti-cheat oracle calculus (U1--U6) with 52
oracles and 52 paired adversaries, paper-grounded falsification against 52
fixture files with bibliographic provenance, the OCEL/OCPN process-truth
machinery, the BLAKE3 receipt chain, and the release certificate evidencing
$60/60$ algorithm reachability and $60$ positive, $120$ negative, and $60$
invariant behavior cases. Every numeric claim in this document was verified
against the repository at the \projver{} working tree before being written
down.
\end{abstract}

\tableofcontents

%=====================================================================
\chapter{Introduction}
\label{ch:intro}

\section{What wasm4pm is}

wasm4pm is a process-mining platform. Its purpose is the classical
process-mining triad --- \emph{discovery} of process models from event logs,
\emph{conformance} checking of logs against models, and \emph{enhancement}
of models with performance and predictive information
\cite{vdaalst2016} --- delivered as a deterministic WebAssembly core that
runs identically in Node.js, browsers, and embedded hosts.

The repository at release \projver{} contains two principal layers and a
family of supporting crates:

\begin{enumerate}
  \item \textbf{The Rust/WASM core} (\texttt{wasm4pm/}). Sixty
  process-mining algorithms compiled via \texttt{wasm-pack}, exported
  through \texttt{wasm-bindgen}. The API surface, documented in
  \texttt{WASM\_API.md}, comprises 335 public Rust$\to$JavaScript exports
  across more than one hundred source modules, spanning core discovery
  (DFG, Alpha++, Heuristic, Inductive miners; metaheuristic miners using
  ACO, PSO, genetic search, A*, hill climbing, simulated annealing, ILP;
  DECLARE; streaming variants), conformance and alignment, POWL
  partial-order workflow handling, temporal-profile analysis, social
  networks, simulation, ML algorithms, and autonomic control loops.

  \item \textbf{The TypeScript monorepo} (\texttt{packages/} and
  \texttt{apps/}). Twelve workspace packages ---
  \texttt{agents}, \texttt{cognition}, \texttt{config},
  \texttt{contracts}, \texttt{engine}, \texttt{kernel}, \texttt{ml},
  \texttt{observability}, \texttt{planner}, \texttt{supabase},
  \texttt{testing}, plus the applications --- with the published CLI,
  \wpm{}, shipping from \texttt{apps/wasm4pm/} as
  \texttt{@wasm4pm/cli}. A web playground lives at
  \texttt{apps/playground-web/}. The workspace version at the time of
  writing is \texttt{26.6.10} (CalVer \texttt{vYY.M.D}).

  \item \textbf{Supporting crates} (\texttt{crates/}):
  \texttt{prolog8}, a byte-capped deterministic proof engine;
  \texttt{miniml-core}, a SIMD-accelerated machine-learning core;
  \texttt{ocpq}, object-centric process querying; \texttt{wasm4pm-cli},
  a development-only Rust CLI; and \texttt{wasm4pm-cognition}, the
  cognition layer that is the subject of
  Chapters~\ref{ch:periodic}--\ref{ch:papers}.
\end{enumerate}

\section{The doctrine: process truth}

The platform's constitution is borrowed from van der Aalst's program of
evidence-based process science \cite{vdaalst2016,ocel2}:

\begin{quote}
\itshape If the code says it worked but the event log cannot prove a lawful
process happened, then it did not work.
\end{quote}

This is not a slogan; it is implemented as a runtime gate. Every cognition
breed execution derives an object-centric event log (OCEL) from its own
inference trace and replays that log against a declared lifecycle model;
non-conformance (fitness $< 1.0$) converts an apparently successful run into
a hard error (Section~\ref{sec:lifecycle}). Every \wpm{} run emits a BLAKE3
receipt with non-empty input and output hashes. Releases are certified by a
machine-checked certificate (Chapter~\ref{ch:truth}).

\section{Reading guide}

Chapter~\ref{ch:arch} describes the platform architecture.
Chapter~\ref{ch:periodic} presents the cognition periodic table: 52
\textsc{partial\_alive} breeds, the lawful \texttt{run\_breed} lifecycle, and
the ten-rung standing ladder. Chapter~\ref{ch:anticheat} develops the
anti-cheat oracle calculus U1--U6 with its paired adversaries.
Chapter~\ref{ch:papers} documents paper-grounded falsification: every breed
is tested against a fixture derived from its source publication, and the
oracle must both \emph{confirm} the published behavior and \emph{refute} a
deliberately corrupted variant. Chapter~\ref{ch:truth} covers the
process-truth machinery --- OCPN lifecycle models, measured-fitness reports,
receipts, and the release certificate. Chapter~\ref{ch:quality} summarizes
the quality evidence, and Chapter~\ref{ch:conclusion} concludes.

%=====================================================================
\chapter{Platform Architecture}
\label{ch:arch}

\section{The Rust/WASM core}

The mining core lives in \texttt{wasm4pm/} and compiles to a single WASM
binary via \texttt{wasm-pack build --target nodejs}. Its design principles,
stated in \texttt{WASM\_API.md}, treat the WASM boundary as a
\emph{sovereign sanctuary}:

\begin{enumerate}
  \item \textbf{Memory isolation.} Every call executes within linear
  memory, isolated from host non-determinism.
  \item \textbf{Deterministic calculus.} Every algorithm must satisfy
  $f(\mathit{BinaryHash}, \mathit{InputHash}, \mathit{Seed}) \to
  (\mathit{Output}, \mathit{Receipt})$, with bit-exact identity required
  for Rank-1 oracles.
  \item \textbf{Immutable verdicts.} A \texttt{Refused} verdict from the
  kernel is process law; no administrative layer can overturn it.
\end{enumerate}

The exported module families (counts as documented and verified in
\texttt{WASM\_API.md} at the current tree) are summarized in
Table~\ref{tab:wasmmodules}.

\begin{table}[ht]
\centering
\caption{WASM core module families (from \texttt{WASM\_API.md}, 335 total
exports).}
\label{tab:wasmmodules}
\begin{tabular}{lrl}
\toprule
Module family & Exports & Purpose \\
\midrule
Core Discovery & 15 & DFG, Alpha++, Heuristic, Inductive, ACO, PSO, \\
 & & genetic, A*, hill climbing, simulated annealing, \\
 & & ILP, DECLARE, streaming variants \\
ML Algorithms & 6+ & classification, clustering, forecasting, anomaly \\
 & & detection, regression, PCA \\
Streaming & 20+ & real-time DFG, SIMD acceleration, conformance \\
POWL & 8+ & partial-order workflow parse/simplify/convert \\
Analysis & 20+ & conformance, simulation, performance, temporal, \\
 & & social networks \\
Prolog8 & 3 & query evaluation, proof verification, capabilities \\
TrueX & 1 & OCEL~2.0 execution-receipt verification \\
 & & (BLAKE3 + JCS-OCEL canonicalization) \\
Autonomic & 10+ & RL orchestrator, SPC, circuit breaker, self-healing \\
AutoMembrane & 25+ & motion classification, actor/object/route/ML/time \\
 & & envelopes, benchmarking \\
\bottomrule
\end{tabular}
\end{table}

Two deliberately documented API quirks illustrate the project's
verification culture: \texttt{discover\_astar} returns
\texttt{iterations\_used} (an iteration count, explicitly \emph{not} a
fitness score), and the metaheuristic miners (\texttt{discover\_aco\_algorithm},
\texttt{discover\_pso\_algorithm}) return a typed \texttt{Err("no\_edges")}
rather than fabricating a model when the log has no directly-follows edge
vocabulary. Refusal is a first-class outcome throughout the system.

\section{The TypeScript monorepo and the \wpm{} CLI}

The monorepo (pnpm workspaces, version \texttt{26.6.10}) layers
typed contracts and orchestration over the WASM core:

\begin{itemize}
  \item \texttt{packages/kernel} --- the WASM loader and execution kernel;
  the published npm tarball \texttt{wasm4pm-26.6.10.tgz} is smoke-tested
  from \texttt{packages/kernel/} during release certification.
  \item \texttt{packages/engine}, \texttt{packages/contracts},
  \texttt{packages/config} --- algorithm orchestration, typed input/output
  contracts, and configuration (environment prefix
  \texttt{WASM4PM\_*}; exit codes $0$ ok, $1$ config, $2$ source,
  $3$ exec, $4$ partial, $5$ system).
  \item \texttt{packages/cognition} --- the TypeScript face of the
  cognition layer; field contracts mirror the Rust WASM output exactly
  (e.g.\ \texttt{cognition\_run} returns
  \texttt{\{status, breed, run\_id, output\_hash, replay\_pointer,}
  \texttt{options\_profile, output\}}).
  \item \texttt{packages/observability} --- OpenTelemetry span emission;
  project rule: every public operation emits spans with
  \texttt{service\_name} and a \texttt{status} of \texttt{"ok"} or
  \texttt{"error"}, and no success claim is accepted without span
  evidence.
  \item \texttt{packages/agents}, \texttt{packages/planner},
  \texttt{packages/ml}, \texttt{packages/supabase},
  \texttt{packages/testing} --- agent harnesses, planning, ML bindings,
  persistence, and the shared test harness.
\end{itemize}

The published CLI, \wpm{} (\texttt{apps/wasm4pm/}, npm package
\texttt{@wasm4pm/cli}), exposes mining and cognition commands, including
\texttt{wpm run}, \texttt{wpm cognition run}, and
\texttt{wpm compile --spec <f.json> [--run]}, a multi-stage compilation
pipeline in which an unknown breed exits with code~2. Every
\texttt{wpm run} and \texttt{wpm cognition run} writes a BLAKE3 receipt to
\texttt{.wasm4pm/receipts/latest.json}; there is deliberately no
\texttt{--no-receipt} escape hatch.

\section{Supporting crates}

\subsection{prolog8 --- byte-capped proof engine}

\texttt{crates/prolog8} is a compact, deterministic inference engine that
evaluates queries against fact blocks and Horn rules, emitting typed proof
DAGs with BLAKE3 receipt chains. It is designed for bounded-memory embedded
use: hard byte caps (e.g.\ a maximum predicate arity of 8) make worst-case
resource consumption a static property. Its use cases inside the platform
are proof-of-admissibility gates in autonomic subsystems and receipt-chain
validation in the manufacturing-style pipelines. Source modules:
\texttt{admission.rs}, \texttt{catalog.rs}, \texttt{hash.rs},
\texttt{kernel.rs}, \texttt{replay.rs}, \texttt{types.rs}, \texttt{wasm.rs}.

\subsection{miniml-core --- SIMD machine learning}

\texttt{crates/miniml-core} (published as \texttt{wminml}) is a minimal
machine-learning core: 70+ algorithms across 15 families, SIMD-accelerated
matrix operations using \texttt{v128} intrinsics, zero-copy serde
serialization, and AutoML via genetic-algorithm feature selection with PSO
hyperparameter optimization. It backs the ML exports of the WASM core.

\subsection{ocpq --- object-centric process querying}

\texttt{crates/ocpq} implements object-centric process querying in the
tradition of K\"usters and van der Aalst's OCPQ line of work
\cite{ocpq}: binding-based queries over OCEL~2.0 logs in which constraints
range over events \emph{and} objects simultaneously, rather than flattening
the log to a single case notion. The crate ships with its own test and
benchmark suites (\texttt{crates/ocpq/tests}, \texttt{crates/ocpq/benches}).

\subsection{wasm4pm-cli and wasm4pm-cognition}

\texttt{crates/wasm4pm-cli} is a development-only Rust CLI (the shipped CLI
is the TypeScript \wpm{}). \texttt{crates/wasm4pm-cognition} is the
cognition layer, treated at length in the remainder of this monograph; it
builds both natively and for \texttt{wasm32-unknown-unknown} (the wasm32
gate is a required check), and exposes \texttt{cognition\_run},
\texttt{cognition\_verify}, and \texttt{system\_build} through
\texttt{wasm-bindgen}.

%=====================================================================
\chapter{The Cognition Periodic Table}
\label{ch:periodic}

\section{Breeds as elements}

The cognition layer organizes classical symbolic-AI reasoning systems as
\emph{breeds} --- each a faithful, deterministic reconstruction of one
historical reasoning paradigm, normalized behind one trait
(\texttt{CognitionBreed}) and one WASM export (\texttt{cognition\_run}).
The metaphor is a periodic table: each breed is an element with a fixed
identity, a family, a historical ancestor, and a measured standing.

The registry, \texttt{crates/wasm4pm-cognition/breeds/registry.json},
contains \textbf{55 entries} at \projver{}, of which \textbf{52 carry
status \textsc{partial\_alive}} and 3 remain \textsc{unsupported}. Each
entry binds a breed to its schemas, oracle suite, OCEL model, and receipt
schema. A representative entry:

\begin{lstlisting}[language={}]
{
  "breed_id": "mycin",
  "breed_name": "MYCIN",
  "historical_ancestor": "MYCIN (Shortliffe, 1976)",
  "generalized_family": "rule_based_cf",
  "input_schema": "mycin_input_v1",
  "output_schema": "mycin_output_v1",
  "specification_relation": "forward_chaining_cf",
  "oracle_suite_id": "oracle_mycin_v1",
  "ocel_model_id": "ocel_mycin_v1",
  "receipt_schema_id": "receipt_mycin_v1",
  "wasm_export_name": "cognition_run",
  "status": "PARTIAL_ALIVE",
  "standing": "DISPATCHABLE"
}
\end{lstlisting}

The standings recorded in the registry at \projver{} are: \textbf{36
\textsc{dispatchable}}, \textbf{16 \textsc{traceable}}, and \textbf{3
\textsc{registered}} (the latter being the three \textsc{unsupported}
entries). The breed source tree, \texttt{src/breeds/}, holds the 52
implementations plus dispatch, standing, and support modules.

The families span the breadth of symbolic AI: rule-based reasoning with
certainty factors (MYCIN), heuristic search and means--ends analysis (GPS,
STRIPS), logic programming and its abductive, inductive, constraint, and
probabilistic extensions (Prolog, abductive LP, ILP, CLP, ProbLog), answer
set programming, non-monotonic logics (default logic, circumscription),
description logic ($\mathcal{EL}$ subsumption), temporal reasoning (Allen's
interval algebra, event calculus, situation calculus, LTL monitoring, CTL
model checking), uncertainty calculi (Bayesian networks, Dempster--Shafer,
fuzzy logic, Markov logic), planning (HTN, partial-order, contingent, MDP,
POMDP), learning (version spaces, EBL, Q-learning, case-based reasoning),
cognitive architectures (Soar, ACT-R, episodic memory), language and
knowledge structures (ELIZA, conceptual-dependency scripts, construction
grammar, frames), qualitative and naive physics, analogy (SME), blackboard
architectures (Hearsay-II), SAT solving (CDCL), tableaux proof, belief
merging, and meta-reasoning.

\section{The lawful lifecycle: \texttt{run\_breed}}
\label{sec:lifecycle}

No breed is invoked directly. All execution flows through
\texttt{run\_breed} in \texttt{src/breeds/dispatch.rs}, which enforces a
five-stage lawful lifecycle:

\begin{enumerate}
  \item \textbf{Preconditions.} \texttt{b.preconditions(input)} validates
  the contract and the breed's declared complexity caps; failure is a
  refusal, not a degraded run.
  \item \textbf{Run.} \texttt{b.run(input)} executes the algorithm.
  \item \textbf{FM-5 gate.} A framework-level fraud gate, independent of
  the breed's own checks: an empty \texttt{inference\_trace} on a
  ``successful'' run is rejected as an FM-5 fraud signal. A breed cannot
  claim success without showing its work.
  \item \textbf{Postconditions.} \texttt{b.postconditions(input, \&output)}
  verifies output-side invariants.
  \item \textbf{OCEL derivation and conformance.} The inference trace is
  converted into an object-centric event log
  (\texttt{derive\_ocel}); if a lifecycle model is declared for the
  breed, the log is replayed against it
  (\texttt{validate\_ocel\_alignment}) and any non-conformance aborts
  the run with the measured fitness in the error message.
\end{enumerate}

The decisive excerpt (verbatim from \texttt{dispatch.rs}):

\begin{lstlisting}[language=Rust]
pub fn run_breed(b: &dyn CognitionBreed, input: &BreedInput)
    -> Result<BreedOutput, String>
{
    b.preconditions(input)
        .map_err(|e| format!("{}: precondition failed: {}", b.id(), e))?;
    let mut output = b.run(input)
        .map_err(|e| format!("{}: {}", e.breed, e.message))?;
    // Framework-level FM-5 gate: a breed cannot claim success without an
    // inference trace, independent of its own postconditions.
    if output.inference_trace.is_empty() {
        return Err(format!(
            "{}: empty inference trace (FM-5 fraud signal, framework gate)",
            b.id()));
    }
    b.postconditions(input, &output)
        .map_err(|e| format!("{}: postcondition failed: {}", b.id(), e))?;

    // Derive OCEL and validate conformance (van der Aalst doctrine)
    let breed_id = format!("{}", b.id());
    let trace_str =
        serde_json::to_string(&output.inference_trace).unwrap_or_default();
    let tmp_run_id =
        blake3::hash(trace_str.as_bytes()).to_hex().to_string();
    let ocel_log =
        crate::ocel::derive_ocel(&breed_id, &tmp_run_id,
                                 &output.inference_trace);
    if let Some(model) = crate::ocel::lifecycle_model_for(&breed_id) {
        let conformance =
            crate::ocel::validate_ocel_alignment(&ocel_log, model);
        if !conformance.is_conforming {
            return Err(format!(
                "{}: OCEL conformance failure (fitness={:.3}): {}",
                breed_id, conformance.fitness,
                conformance.refusals.join("; ")));
        }
    }
    output.ocel_log =
        Some(serde_json::to_value(&ocel_log).unwrap_or(
             serde_json::Value::Null));
    Ok(output)
}
\end{lstlisting}

Dispatch itself (\texttt{dispatch\_breed}) parses the breed string into the
\texttt{BreedId} enum and matches \emph{exhaustively}: the Rust compiler
enforces that every named breed has a dispatch arm, so a breed cannot be
registered yet silently unreachable.

Determinism is structural: all breed collections are
\texttt{BTreeMap}/\texttt{BTreeSet} or sorted vectors (hash-map iteration
order would break determinism and receipts); the only permitted RNG is the
seeded \texttt{support::rng::seeded\_rng()}; and a shared
\texttt{assert\_deterministic} harness compares full \texttt{BreedOutput}
bytes across repeated runs (hand-rolled determinism checks are themselves
classified as a cheat, code A11).

\section{The ten-rung standing ladder}

A breed's maturity is not a boolean. \texttt{src/breeds/standing.rs}
defines \texttt{BreedStanding}, a strictly ordered ten-rung ladder; the
documentation comments on the enum variants are reproduced in
Table~\ref{tab:ladder}.

\begin{table}[ht]
\centering
\caption{The \texttt{BreedStanding} ladder (\texttt{src/breeds/standing.rs}).}
\label{tab:ladder}
\begin{tabular}{clp{8.6cm}}
\toprule
Rung & Standing & Requirement \\
\midrule
1 & \textsc{Named} & \texttt{BreedId} variant exists in the enum and \texttt{BreedId::ALL}. \\
2 & \textsc{Registered} & Entry present in \texttt{breeds/registry.json} with mandatory fields. \\
3 & \textsc{Dispatchable} & Dispatch arm exists in \texttt{dispatch\_breed\_id} (compile-enforced). Minimum standing for \textsc{partial\_alive} status. \\
4 & \textsc{Bounded} & \texttt{complexity\_caps} declared; \texttt{preconditions()} enforces them. \\
5 & \textsc{Oracled} & Oracle fixture exists; paper \texttt{expected\_value} asserted. \\
6 & \textsc{Traceable} & \texttt{inference\_trace} non-empty; trace-law checks pass. \\
7 & \textsc{Canonical} & BLAKE3 receipt deterministic; \texttt{SortedFacts} used throughout. \\
8 & \textsc{Refusable} & \texttt{CognitionError} variants in use; all refusal paths tested. \\
9 & \textsc{Replayable} & OCEL conformance fitness $= 1.0$; \texttt{replay\_pointer} valid. \\
10 & \textsc{Certified} & Ed25519 signature verified; Merkle root anchored; counter-test green. \\
\bottomrule
\end{tabular}
\end{table}

The constant \texttt{BreedStanding::PARTIAL\_ALIVE\_MINIMUM} is
\textsc{Dispatchable}: a registry entry may flip from
\textsc{unsupported} to \textsc{partial\_alive} only after the build and
test gates are green \emph{and} an OCEL report with measured-fitness
provenance exists for the breed. At \projver{} the population sits at
rungs~3--6: 36 breeds at \textsc{dispatchable}, 16 at \textsc{traceable},
3 (unsupported) at \textsc{registered}. The ladder makes ``how alive is
this breed?'' a queryable, monotone fact rather than a vibe.

\section{Contracts at the WASM boundary}

The cognition layer's WASM surface (\texttt{src/wasm.rs}) is contract-exact
and mirrored verbatim into the TypeScript layer:

\begin{itemize}
  \item \texttt{cognition\_run} input:
  \verb|{ breed, contract, options?: { profile? } }| ---
  serde \texttt{deny\_unknown\_fields} rejects bare breed inputs.
  \item \texttt{cognition\_run} output (\texttt{ContractResult}):
  \verb|{ status: "ok", breed, run_id, output_hash,|
  \verb|replay_pointer, options_profile, output }|, where
  \texttt{replay\_pointer} is the first 16 hex characters of
  \texttt{output\_hash}.
  \item \texttt{cognition\_verify} output (\texttt{VerifyResult}):
  \verb+{ findings, status: "verified" | "has_findings" }+.
  \item \texttt{system\_build} output:
  \verb|{ pareto_front, dominated }| --- breed-system composition as
  multi-objective selection.
\end{itemize}

Repository hooks (\texttt{cognition-contract-guard.sh},
\texttt{cognition-wasm-gate.sh}) enforce that no TypeScript consumer
invents fields the Rust side does not emit.

%=====================================================================
\chapter{The Anti-Cheat Oracle Calculus}
\label{ch:anticheat}

\section{Threat model}

The cognition layer assumes its own implementations are adversarial. The
threat model (\texttt{docs/breeds/anti-cheat-threat-model.md}) catalogs the
ways a breed could \emph{appear} to implement a classical algorithm while
actually doing something cheaper: hard-coding answers for known fixtures,
emitting hollow traces, ignoring its input, never refusing malformed input,
or faking determinism. Each cheat mode carries an \texttt{AC-*} taxonomy
code (A1--A11), and each breed is required to have a \emph{counter-test}
that would catch its most tempting cheat. A representative entry: for
\texttt{csp\_ac3}, the feared cheat is ``a backtracking-only solver that
never does arc revision (gets right answers, fakes \texttt{csp-revise}
steps)''; the counter-test asserts that the $K_3$/2-coloring UNSAT
instance is refuted by AC-3 \emph{alone} --- the trace must show a
domain-wipeout \texttt{csp-revise} step and \emph{zero} \texttt{csp-assign}
steps, cross-checked by a CLP instance that must solve with zero
backtracks.

\section{The universal laws U1--U6}

\texttt{src/breeds/support/oracle.rs} defines the universal anti-cheat
calculus. \texttt{run\_universal\_anticheat} executes, in order, U1, U2,
U2b, U3, U4, U5 and returns exactly six results; all must pass for a breed
to be \emph{oracle-green}. The mapping from law to cheat mode, taken from
the module documentation, is given in Table~\ref{tab:ulaws}.

\begin{table}[ht]
\centering
\caption{Universal anti-cheat laws (\texttt{support/oracle.rs}).}
\label{tab:ulaws}
\begin{tabular}{clp{7.6cm}}
\toprule
Law & Targets & Check \\
\midrule
U1 & A1, A2 & A \emph{novel} input (not present in any fixture) runs without error --- defeats fixture memorization. \\
U2 & A4 & \texttt{assert\_intermediate} passes on the novel trace --- the trace must contain the algorithm's real intermediate steps. \\
U2b & A4 & \texttt{assert\_trace\_values} passes (Surface-3 content check) --- step \emph{values}, not just step kinds, must be meaningful. \\
U3 & A3, A6 & A boundary pair of inputs must yield different outputs --- defeats constant-output and input-insensitive breeds. \\
U4 & A5 & A designated refusal input must return \texttt{Err} --- a breed that never refuses is cheating. \\
U5 & A11 & Determinism over full \texttt{BreedOutput} bytes --- defeats sham determinism. \\
U6 & meta & Adversary check: a known-cheating implementation must \emph{fail} the oracle --- the oracle itself is tested. \\
\bottomrule
\end{tabular}
\end{table}

If U1 fails, U2 and U2b are recorded as failed with the explicit detail
``skipped because U1 failed'' --- skipping is never silent.

\section{52 oracles, 52 adversaries}

Each of the 52 breeds has a per-breed oracle implementation under
\texttt{src/breeds/support/oracle\_impls/} (organized by family:
\texttt{dialogue.rs}, \texttt{learning.rs}, \texttt{logic.rs},
\texttt{planning.rs}, \texttt{rule\_fact.rs}). Each oracle supplies the
novel input, intermediate-step assertions, boundary pair, and refusal input
for its breed --- and, crucially, each oracle is paired with an
\emph{adversary}: a deliberately cheating implementation carrying its
\texttt{AC-*} code, used by U6 to verify the oracle has discriminating
power. The framework ships a generic hollow-trace adversary (cheat code
\texttt{AC-A4}) that emits a trace of meaningless kinds; the meta-oracle
asserts not merely that the adversary fails but that the failure detail
\emph{names the cheat code}:

\begin{lstlisting}[language=Rust]
assert!(detail.contains("AC-A4"),
        "detail must name cheat code: {}", detail);
\end{lstlisting}

This closes the classic self-reference loophole (FM-5): a test suite that
only ever runs honest implementations can never demonstrate that it would
catch a dishonest one. Here, for every breed, the suite demonstrably
rejects a known cheat.

\section{Fresh-name hygiene}

Hidden-oracle inputs use \texttt{uo\_*}-prefixed fresh names (universal
oracle namespace). A grep gate in \texttt{anti\_fraud\_gate.rs} enforces
that no \texttt{uo\_*} fresh name appears anywhere in \texttt{src/breeds/}
--- a breed implementation that special-cases the oracle's secret inputs
(cheat A8) is caught at test time by textual quarantine, independent of any
behavioral check.

\section{The integration suite}

\texttt{tests/universal\_anticheat.rs} contains \textbf{54 test
functions} (verified by counting \verb|#[test]| attributes), covering all
52 breeds plus framework-level meta-tests. All 54 pass at \projver{}.

%=====================================================================
\chapter{Paper-Grounded Falsification}
\label{ch:papers}

\section{Method: confirm and refute}

Every breed must be anchored to its source publication, not to the
implementer's intuition. \texttt{tests/paper\_falsification.rs} iterates
over the fixture directory \texttt{tests/fixtures/papers/} --- \textbf{52
fixture files}, one per \textsc{partial\_alive} breed --- and for each
fixture performs both directions of a Popperian check:

\begin{enumerate}
  \item \textbf{Confirm}: dispatch the fixture input through the full
  lawful path (\texttt{run\_breed}, i.e.\ including the FM-5 gate and OCEL
  conformance) and assert the published expected value.
  \item \textbf{Refute}: dispatch a corrupted variant and assert the
  oracle rejects it --- a fixture that cannot fail is not evidence.
\end{enumerate}

Each fixture carries a \texttt{provenance} block with a full bibliographic
citation (and, where the project archive holds the PDF, a path to it), and
paper-derived fixtures must assert the \emph{published} number --- for
example, the Bayesian-network fixture asserts Pearl's value
$0.284171835$ to full precision rather than a rounded approximation, and
the ProbLog fixture asserts the success probability $0.552$.

\section{Breeds, source papers, and distinguishing assertions}

Table~\ref{tab:papers} samples 18 representative breeds. Citations are
transcribed from the on-disk fixture provenance blocks; the
``distinguishing assertion'' column states the fixture's
\texttt{expected.value} or the behavioral property the oracle pins down.

\begin{small}
\begin{longtable}{p{2.5cm}p{6.2cm}p{4.9cm}}
\caption{Paper-grounded fixtures (sample of 18 of 52;
\texttt{tests/fixtures/papers/*.json}).}
\label{tab:papers}\\
\toprule
Breed & Source paper (fixture provenance) & Distinguishing assertion \\
\midrule
\endfirsthead
\toprule
Breed & Source paper (fixture provenance) & Distinguishing assertion \\
\midrule
\endhead
\bottomrule
\endfoot
mycin & Shortliffe \& Buchanan (1975), \emph{A model of inexact reasoning in medicine} \cite{shortliffe1975} & Certainty-factor accumulation across chained rules reproduces the CF-model combination values. \\
strips & Fikes \& Nilsson (1971), \emph{STRIPS: A new approach to the application of theorem proving to problem solving} \cite{fikes1971} & Plan found via add/delete-list operator application on the published blocks domain. \\
gps & Newell \& Simon (1963), \emph{GPS, a program that simulates human thought} \cite{newell1963} & Means--ends analysis reduces the published difference table to a solved goal stack. \\
eliza & Weizenbaum (1966), \emph{ELIZA --- a computer program for the study of natural language communication between man and machine} \cite{weizenbaum1966} & Keyword-rank decomposition/reassembly reproduces script-driven responses. \\
bayesian\_network & Pearl (1988), \emph{Probabilistic Reasoning in Intelligent Systems} \cite{pearl1988} & Posterior asserted to published precision: $0.284171835$. \\
dempster\_shafer & Shafer (1976), \emph{A Mathematical Theory of Evidence} \cite{shafer1976} & Combined belief $0.99$ via Dempster's rule on the fixture mass functions. \\
fuzzy\_logic & Mamdani \& Assilian (1975), \emph{An experiment in linguistic synthesis with a fuzzy logic controller} \cite{mamdani1975} & Defuzzified output $41.66667$ (centroid of clipped consequents). \\
default\_logic & Reiter (1980), \emph{A logic for default reasoning} \cite{reiter1980} & Extension computation on the canonical default theory. \\
circumscription & McCarthy (1980), \emph{Circumscription --- a form of non-monotonic reasoning} \cite{mccarthy1980} & Minimal-model result: \texttt{flies\_tweety=true}, \texttt{flies\_opus=false}. \\
allen\_temporal & Allen (1983), \emph{Maintaining knowledge about temporal intervals} \cite{allen1983} & Path-consistency over the 13 interval relations yields \texttt{temporal-consistent}. \\
csp\_ac3 & Mackworth (1977), \emph{Consistency in networks of relations} \cite{mackworth1977} & $K_3$/2-coloring refuted by arc revision alone: a \texttt{csp-revise} domain wipeout and zero \texttt{csp-assign} steps. \\
asp & Gelfond \& Lifschitz (1988), \emph{The stable model semantics for logic programming} \cite{gelfond1988} & Stable models of the published program enumerated exactly. \\
clp & Jaffar \& Lassez (1987), \emph{Constraint logic programming} \cite{jaffar1987} & Solution \texttt{x=6, y=3} by propagation (zero backtracks on the pure-propagation instance). \\
version\_space & Mitchell (1982), \emph{Generalization as search} \cite{mitchell1982} & S/G boundary-set convergence on the published example sequence. \\
rl\_symbolic & Watkins \& Dayan (1992), \emph{Q-learning} \cite{watkins1992} & Q-value convergence verified on the fixture MDP. \\
analogy\_sme & Falkenhainer, Forbus \& Gentner (1989), \emph{The Structure-Mapping Engine} \cite{falkenhainer1989} & Mapping \texttt{sun$\to$nucleus}, \texttt{planet$\to$electron} for the solar-system/atom analogy. \\
sat\_cdcl & Marques-Silva \& Sakallah (1999), \emph{GRASP: a search algorithm for propositional satisfiability} \cite{marquessilva1999} & \textsc{unsat} verdict with conflict-driven clause learning on the fixture formula. \\
problog & De Raedt, Kimmig \& Toivonen (2007), \emph{ProbLog: a probabilistic Prolog and its application in link discovery} \cite{deraedt2007} & Success probability $0.552000$ on the fixture program. \\
\end{longtable}
\end{small}

The remaining 34 fixtures follow the same pattern, citing inter alia
Anderson \& Lebiere (ACT-R), Laird, Newell \& Rosenbloom (Soar)
\cite{laird1987}, Kowalski (predicate logic as programming language)
\cite{kowalski1974}, Kakas, Kowalski \& Toni (abductive LP), Quinlan
(FOIL), Minsky (frames) \cite{minsky1974}, Schank \& Abelson (scripts),
Erman et al.\ (Hearsay-II), Feigenbaum, Buchanan \& Lederberg (DENDRAL),
Aamodt \& Plaza (CBR), Clarke, Emerson \& Sistla (CTL), Havelund \& Rosu
(LTL monitoring), Nau et al.\ (SHOP2 HTN), McAllester \& Rosenblitt
(systematic nonlinear planning), Kaelbling, Littman \& Cassandra (POMDP),
Bellman (dynamic programming), Richardson \& Domingos (Markov logic),
Konieczny \& Pino P\'erez (belief merging), Baader, Brandt \& Lutz
($\mathcal{EL}$), Smullyan (tableaux), de Kleer \& Brown (qualitative
physics), Hayes (naive physics), Tulving and Nuxoll \& Laird (episodic
memory), Goldberg (construction grammar), Marr \& Poggio (stereo
disparity), Sussman (HACKER), and Cox \& Raja (metareasoning). Several
fixtures are marked \texttt{citation-only} where no PDF is archived; in
those cases the citation alone is the provenance and no quoted text is
asserted.

%=====================================================================
\chapter{Process Truth: OCEL, Receipts, and the Release Certificate}
\label{ch:truth}

\section{OCPN lifecycle models}

Each breed declares an L1 lifecycle model as an object-centric Petri net
in \texttt{ocel/models/l1/<breed\_id>.ocpn.json} --- \textbf{52 models},
one per \textsc{partial\_alive} breed (verified by directory count). The
model names the lawful step kinds and their ordering; for MYCIN, for
example, the declared steps are \texttt{rule\_selected},
\texttt{premise\_matched}, \texttt{cf\_accumulated},
\texttt{threshold\_checked}, \texttt{diagnosis\_emitted}.

\section{Measured-fitness reports}

\texttt{ocel/reports/} holds \textbf{52 fitness reports}, one per breed.
Each records the model ID, model level, report date, method, and the
measured fitness with provenance. The MYCIN report (verbatim fields):

\begin{lstlisting}[language={}]
{
  "breed_id": "mycin",
  "model_id": "mycin-l1-v1",
  "model_level": "L1",
  "report_date": "2026-06-10",
  "method": "native_dfa_replay",
  "fitness": 1,
  "fitness_note": "Alignment fitness = 1.0 measured by
    validate_ocel_alignment() (native Rust DFA replay).
    Conformance gate enforced in run_breed()."
}
\end{lstlisting}

Project law: fitness above $0.85$ is required for a model to be considered
valid at all, and route admission into the broader MCPP ecosystem requires
fitness of exactly $1.0$. Because the conformance check runs \emph{inside}
\texttt{run\_breed}, the reports are not aspirational documentation --- a
breed whose traces stopped conforming would start failing at runtime.

\section{The BLAKE3 receipt chain}

Every \texttt{wpm run} and \texttt{wpm cognition run} emits a receipt to
\texttt{.wasm4pm/receipts/latest.json} with non-empty \texttt{input\_hash}
and \texttt{output\_hash} (BLAKE3). On the Rust side, the
\texttt{run\_id} is itself a BLAKE3 hash, and \texttt{replay\_pointer} (the
first 16 hex characters of \texttt{output\_hash}) lets any consumer
re-locate and re-verify the exact output later. The TrueX module in the
WASM core verifies OCEL~2.0 execution receipts using BLAKE3 over a
JCS-canonicalized OCEL profile, so receipt verification is available inside
the sandbox, not only in the host.

\section{The release certificate}

\texttt{RELEASE\_CERTIFICATE.v26.6.10.json} is the machine-checked release
artifact for \projver{}. Its verified contents are summarized in
Table~\ref{tab:cert}.

\begin{table}[ht]
\centering
\caption{Release certificate \projver{}
(\texttt{RELEASE\_CERTIFICATE.v26.6.10.json}).}
\label{tab:cert}
\begin{tabular}{lp{8.6cm}}
\toprule
Field & Value \\
\midrule
Package & \texttt{wasm4pm} 26.6.10, git commit \texttt{eceacfcb43d5\ldots} \\
Reachability & 60 of 60 algorithms reachable; reachability hash \texttt{92e0264e\ldots} \\
Behavior & 60 algorithms $\times$ (60 positive, 120 negative, 60 invariant) cases; \texttt{all\_failed\_correctly: true}; evidence hash \texttt{b79d5e7f\ldots} \\
Examples & 8 examples, 64 total executions; manifest hash \texttt{86933a9b\ldots} \\
Artifact & \texttt{wasm4pm-26.6.10.tgz}; sha512 integrity \texttt{A6g+gOqY\ldots}; pack-smoke path \texttt{packages/kernel/}; WASM bundle hash \texttt{d4bc2f40\ldots} \\
\bottomrule
\end{tabular}
\end{table}

Three properties deserve emphasis. First, \emph{reachability} is asserted
per algorithm: all 60 mining algorithms are demonstrably callable through
the published package, not merely compiled in. Second, the behavior
evidence is dominated by \emph{negative} cases --- 120 negative against 60
positive --- and the certificate's \texttt{all\_failed\_correctly: true}
asserts that every injected invalid input was refused with the correct
error, in line with the negative-testing requirement of the process-truth
constitution. Third, the certificate pins the exact tarball and WASM
binary by hash, closing the gap between ``the tests passed'' and ``the
artifact you installed is the artifact that passed.''

%=====================================================================
\chapter{Quality Evidence}
\label{ch:quality}

\section{Test suites at \projver{}}

The evidence base, with each figure verified against the working tree:

\begin{itemize}
  \item \textbf{Universal anti-cheat}: 54 test functions in
  \texttt{tests/universal\_anticheat.rs} (counted directly), all green ---
  every breed passes U1--U5 + U2b, and every paired adversary is rejected
  by U6 with its cheat code named.
  \item \textbf{Paper falsification}: \texttt{tests/paper\_falsification.rs}
  drives all 52 paper fixtures (counted directly in
  \texttt{tests/fixtures/papers/}) through the confirm-and-refute
  protocol via the full lawful dispatch path.
  \item \textbf{Cargo suites}: the cognition crate's unit, integration,
  and anti-fraud suites are green at the release commit; the \projver{}
  integration commit records 868 cargo tests and 237 vitest tests passing
  across the workspace.
  \item \textbf{wasm32 gate}: \texttt{cargo check --target
  wasm32-unknown-unknown --features wasm} is a required release gate; the
  cognition layer builds for the WASM target, and the WASM bundle hash is
  pinned in the release certificate.
  \item \textbf{TypeScript}: the monorepo builds and tests under pnpm at
  workspace version 26.6.10; the FM-5 rule requires at least one
  cognition integration test to exercise the \emph{real} WASM binary with
  no mocking of \texttt{init.js}.
  \item \textbf{Release certification}: $60/60$ reachability,
  $60/120/60$ behavior cases with \texttt{all\_failed\_correctly},
  tarball and bundle integrity hashes (Table~\ref{tab:cert}).
  \item \textbf{OCEL conformance}: 52 measured-fitness reports, each
  recording fitness $1.0$ by native DFA replay, with the same check
  enforced at runtime inside \texttt{run\_breed}.
\end{itemize}

\section{Defense in depth}

The quality argument is layered so that no single artifact has to be
trusted:

\begin{enumerate}
  \item The \emph{compiler} enforces dispatch totality (exhaustive enum
  match) and collection determinism (BTree types).
  \item The \emph{framework} enforces the FM-5 empty-trace gate and OCEL
  conformance on every run, in production code, not only in tests.
  \item The \emph{oracles} enforce behavioral fidelity (U1--U5), and the
  \emph{adversaries} enforce oracle strength (U6).
  \item The \emph{papers} anchor expected values to published numbers
  with bibliographic provenance, and the refute direction proves the
  fixtures can fail.
  \item The \emph{receipts and certificate} anchor the shipped artifact
  to the evidence by cryptographic hash.
\end{enumerate}

A failure anywhere in this stack is, by project law, an Andon event: the
line stops on \texttt{error[E}, \texttt{test.*FAILED}, \texttt{FM-5
violation}, or \texttt{panicked at}, and the defect is fixed before any
other work proceeds.

%=====================================================================
\chapter{Conclusion}
\label{ch:conclusion}

wasm4pm at \projver{} demonstrates that a process-mining platform can hold
itself to the standard it applies to the processes it mines. The mining
core supplies 60 algorithms behind a deterministic WASM boundary; the
cognition layer revives 52 classical reasoning systems as breeds, each
paper-anchored, oracle-guarded, adversary-tested, and conformance-gated at
runtime; and the release machinery converts ``it works'' from an assertion
into a certificate. The discipline throughout is van der Aalst's, applied
reflexively: the system does not ask to be believed --- it asks to be
replayed.

Future work tracked in the repository includes advancing breeds up the
standing ladder toward \textsc{Replayable} and \textsc{Certified}
(Ed25519-signed, Merkle-anchored), promoting the three remaining
\textsc{registered} entries, and validating the mining algorithms against
external real-world XES/OCEL corpora.

%=====================================================================
\appendix
\chapter{Full Breed Roster}
\label{app:roster}

Table~\ref{tab:roster} reproduces the complete registry
(\texttt{crates/wasm4pm-cognition/breeds/registry.json}, 55 entries) at
\projver{}, sorted by breed identifier. The \texttt{historical\_ancestor}
and \texttt{generalized\_family} columns are transcribed verbatim,
including the literal placeholder string \texttt{TBD} where the registry
has not yet been enriched for an entry --- these placeholders are
metadata-completeness debt, not implementation debt; every
\textsc{dispatchable} and \textsc{traceable} breed below has a full
implementation, oracle, adversary, paper fixture, OCPN model, and fitness
report. The three \textsc{registered} entries (\texttt{morphological},
\texttt{ocpm\_route\_discoverer}, \texttt{triz}) are the
\textsc{unsupported} breeds awaiting promotion.

\begin{small}
\begin{longtable}{p{3.4cm}p{5.6cm}p{3.4cm}l}
\caption{Complete breed registry at \projver{} (55 entries; 36
\textsc{dispatchable}, 16 \textsc{traceable}, 3 \textsc{registered}).}
\label{tab:roster}\\
\toprule
Breed & Historical ancestor & Family & Standing \\
\midrule
\endfirsthead
\toprule
Breed & Historical ancestor & Family & Standing \\
\midrule
\endhead
\bottomrule
\endfoot
\texttt{abductive\_ibe} & TBD & TBD & \textsc{traceable} \\
\texttt{abductive\_lp} & Abduction (Peirce, 1878) (abductive) & R\_generalized\_abductive & \textsc{traceable} \\
\texttt{act\_r} & TBD & TBD & \textsc{traceable} \\
\texttt{allen\_temporal} & Temporal Logic (Allen, 1983) (temporal) & R\_generalized\_temporal & \textsc{dispatchable} \\
\texttt{analogy\_sme} & Structure Mapping (Gentner, 1983) (analogical) & R\_generalized\_analogical & \textsc{dispatchable} \\
\texttt{asp} & TBD & TBD & \textsc{traceable} \\
\texttt{autoinstinct\_learning} & autoinstinct\_learning & autoinstinct & \textsc{dispatchable} \\
\texttt{autoinstinct\_neurosis} & autoinstinct\_neurosis & autoinstinct & \textsc{dispatchable} \\
\texttt{autoinstinct\_semantics} & autoinstinct\_semantics & autoinstinct & \textsc{dispatchable} \\
\texttt{autoinstinct\_vision} & autoinstinct\_vision & autoinstinct & \textsc{dispatchable} \\
\texttt{bayesian\_network} & Bayesian Networks (Pearl, 1985) (bayesian) & R\_generalized\_probabilistic & \textsc{traceable} \\
\texttt{belief\_merging} & TBD & TBD & \textsc{traceable} \\
\texttt{cbr} & CBR (Kolodner, 1983) & case\_based\_reasoning & \textsc{dispatchable} \\
\texttt{circumscription} & TBD & TBD & \textsc{traceable} \\
\texttt{clp} & TBD & TBD & \textsc{traceable} \\
\texttt{construction\_grammar} & Construction Grammar (Goldberg, 1995) & grammar\_semantics & \textsc{dispatchable} \\
\texttt{contingent\_plan} & AND-OR belief-state search (Russell \& Norvig, AIMA \S4.3.2) & planning\_under\_uncertainty & \textsc{dispatchable} \\
\texttt{csp\_ac3} & CSP (Mackworth, 1977) (constraint) & R\_generalized\_constraint & \textsc{traceable} \\
\texttt{ctl\_check} & TBD & TBD & \textsc{dispatchable} \\
\texttt{default\_logic} & TBD & TBD & \textsc{dispatchable} \\
\texttt{dempster\_shafer} & Dempster-Shafer (Shafer, 1976) (dempster\_shafer) & R\_generalized\_evidential & \textsc{dispatchable} \\
\texttt{dendral} & DENDRAL (Feigenbaum et al., 1969) & hypothesis\_enumeration & \textsc{dispatchable} \\
\texttt{description\_logic} & Description Logics (Baader et al., 1991) (ontological) & R\_generalized\_ontological & \textsc{dispatchable} \\
\texttt{ebl} & TBD & TBD & \textsc{dispatchable} \\
\texttt{eliza} & ELIZA (Weizenbaum, 1966) & pattern\_matching\_dialogue & \textsc{dispatchable} \\
\texttt{episodic\_memory} & TBD & TBD & \textsc{traceable} \\
\texttt{event\_calculus} & TBD & TBD & \textsc{dispatchable} \\
\texttt{frames\_inheritance} & TBD & TBD & \textsc{dispatchable} \\
\texttt{fuzzy\_logic} & Fuzzy Sets (Zadeh, 1965) (fuzzy) & R\_generalized\_fuzzy & \textsc{dispatchable} \\
\texttt{gps} & GPS (Newell \& Simon, 1957) & means\_ends\_analysis & \textsc{dispatchable} \\
\texttt{hearsay} & Hearsay-II (Erman et al., 1980) & blackboard\_architecture & \textsc{dispatchable} \\
\texttt{htn\_planning} & SHOP2 (Nau et al., 2003) & hierarchical\_task\_network & \textsc{dispatchable} \\
\texttt{ilp} & ILP (Muggleton, 1991) (inductive) & R\_generalized\_inductive & \textsc{dispatchable} \\
\texttt{ltl\_monitor} & LTL Monitor (Havelund \& Ro\c{s}u, 2001) (ltl\_monitor) & R\_generalized\_temporal & \textsc{traceable} \\
\texttt{markov\_logic} & Markov logic networks (Richardson \& Domingos, 2006) & statistical\_relational & \textsc{dispatchable} \\
\texttt{mdp} & TBD & TBD & \textsc{traceable} \\
\texttt{meta\_reasoning} & Metareasoning (Cox \& Raja, 2011) & meta\_level\_control & \textsc{dispatchable} \\
\texttt{morphological} & TBD & TBD & \textsc{registered} \\
\texttt{mycin} & MYCIN (Shortliffe, 1976) & rule\_based\_cf & \textsc{dispatchable} \\
\texttt{naive\_physics} & TBD & TBD & \textsc{dispatchable} \\
\texttt{ocpm\_route\_discoverer} & TBD & TBD & \textsc{registered} \\
\texttt{partial\_order\_plan} & TBD & TBD & \textsc{dispatchable} \\
\texttt{pomdp} & POMDP planning (Kaelbling, Littman \& Cassandra, 1998) & decision\_theoretic & \textsc{traceable} \\
\texttt{problog} & TBD & TBD & \textsc{traceable} \\
\texttt{prolog} & Prolog (Colmerauer \& Roussel, 1972) & logic\_programming & \textsc{dispatchable} \\
\texttt{qualitative\_reason} & TBD & TBD & \textsc{dispatchable} \\
\texttt{rl\_symbolic} & TBD & TBD & \textsc{dispatchable} \\
\texttt{sat\_cdcl} & TBD & TBD & \textsc{traceable} \\
\texttt{script\_sam} & TBD & TBD & \textsc{dispatchable} \\
\texttt{situation\_calculus} & TBD & TBD & \textsc{dispatchable} \\
\texttt{soar} & SOAR (Laird, Newell \& Rosenbloom, 1987) & unified\_cognitive\_architecture & \textsc{dispatchable} \\
\texttt{strips} & STRIPS (Fikes \& Nilsson, 1971) & classical\_planning & \textsc{dispatchable} \\
\texttt{tableaux} & Analytic tableaux (Smullyan, 1968) & proof\_search & \textsc{dispatchable} \\
\texttt{triz} & TBD & TBD & \textsc{registered} \\
\texttt{version\_space} & TBD & TBD & \textsc{traceable} \\
\end{longtable}
\end{small}

%=====================================================================
\begin{thebibliography}{99}

\bibitem{vdaalst2016} W.~M.~P. van der Aalst.
\emph{Process Mining: Data Science in Action}, 2nd edition. Springer, 2016.

\bibitem{ocel2} W.~M.~P. van der Aalst.
\emph{Object-Centric Process Mining: Dealing with Divergence and
Convergence in Event Data}. In: Software Engineering and Formal Methods
(SEFM), 2019.

\bibitem{ocpq} J.~N. K\"usters and W.~M.~P. van der Aalst.
\emph{OCPQ: Object-Centric Process Querying}. (Object-centric process
querying line of work; basis of the \texttt{ocpq} crate.)

\bibitem{shortliffe1975} E.~H. Shortliffe and B.~G. Buchanan.
\emph{A model of inexact reasoning in medicine}. Mathematical Biosciences,
23(3--4), 1975.

\bibitem{fikes1971} R.~E. Fikes and N.~J. Nilsson.
\emph{STRIPS: A new approach to the application of theorem proving to
problem solving}. Artificial Intelligence, 2(3--4), 1971.

\bibitem{newell1963} A.~Newell and H.~A. Simon.
\emph{GPS, a program that simulates human thought}. In: Computers and
Thought, 1963.

\bibitem{weizenbaum1966} J.~Weizenbaum.
\emph{ELIZA --- a computer program for the study of natural language
communication between man and machine}. Communications of the ACM, 9(1),
1966.

\bibitem{pearl1988} J.~Pearl.
\emph{Probabilistic Reasoning in Intelligent Systems: Networks of
Plausible Inference}. Morgan Kaufmann, 1988.

\bibitem{shafer1976} G.~Shafer.
\emph{A Mathematical Theory of Evidence}. Princeton University Press, 1976.

\bibitem{mamdani1975} E.~H. Mamdani and S.~Assilian.
\emph{An experiment in linguistic synthesis with a fuzzy logic controller}.
International Journal of Man-Machine Studies, 7(1), 1975.

\bibitem{reiter1980} R.~Reiter.
\emph{A logic for default reasoning}. Artificial Intelligence, 13(1--2),
1980.

\bibitem{mccarthy1980} J.~McCarthy.
\emph{Circumscription --- a form of non-monotonic reasoning}. Artificial
Intelligence, 13(1--2), 1980.

\bibitem{allen1983} J.~F. Allen.
\emph{Maintaining knowledge about temporal intervals}. Communications of
the ACM, 26(11), 1983.

\bibitem{mackworth1977} A.~K. Mackworth.
\emph{Consistency in networks of relations}. Artificial Intelligence, 8(1),
1977.

\bibitem{gelfond1988} M.~Gelfond and V.~Lifschitz.
\emph{The stable model semantics for logic programming}. In: Proceedings of
the Fifth International Conference and Symposium on Logic Programming,
1988.

\bibitem{jaffar1987} J.~Jaffar and J.-L. Lassez.
\emph{Constraint logic programming}. In: Proceedings of the 14th ACM
SIGACT-SIGPLAN Symposium on Principles of Programming Languages (POPL),
1987.

\bibitem{mitchell1982} T.~M. Mitchell.
\emph{Generalization as search}. Artificial Intelligence, 18(2), 1982.

\bibitem{watkins1992} C.~J.~C.~H. Watkins and P.~Dayan.
\emph{Q-learning}. Machine Learning, 8, 1992.

\bibitem{falkenhainer1989} B.~Falkenhainer, K.~D. Forbus, and D.~Gentner.
\emph{The Structure-Mapping Engine: algorithm and examples}. Artificial
Intelligence, 41(1), 1989.

\bibitem{marquessilva1999} J.~P. Marques-Silva and K.~A. Sakallah.
\emph{GRASP: a search algorithm for propositional satisfiability}. IEEE
Transactions on Computers, 48(5), 1999.

\bibitem{deraedt2007} L.~De Raedt, A.~Kimmig, and H.~Toivonen.
\emph{ProbLog: a probabilistic Prolog and its application in link
discovery}. In: Proceedings of IJCAI, 2007.

\bibitem{laird1987} J.~E. Laird, A.~Newell, and P.~S. Rosenbloom.
\emph{Soar: an architecture for general intelligence}. Artificial
Intelligence, 33(1), 1987.

\bibitem{kowalski1974} R.~A. Kowalski.
\emph{Predicate logic as programming language}. In: IFIP Congress, 1974.

\bibitem{minsky1974} M.~Minsky.
\emph{A framework for representing knowledge}. MIT AI Laboratory Memo 306,
1974.

\end{thebibliography}

\end{document}
